\chapter{Zbieżność topologiczna}
\label{ch:topog-zbieznosc}

\noindent\textbf{Cel rozdziału.}
Zbieżność opisuje, w jaki sposób punkty przestrzeni mogą zbliżać się
w sensie jej topologii. W przestrzeniach metrycznych zwykle posługujemy
się ciągami, ale w dowolnej przestrzeni topologicznej same ciągi
nie muszą wykrywać domknięcia ani ciągłości. Zaczniemy od zbieżności
ciągów i jej związku z pierwszym aksjomatem przeliczalności.
Następnie wprowadzimy
sieci (nazywane też ciągami uogólnionymi), które przywracają
pełne kryteria domknięcia i ciągłości, a na końcu filtry,
będące równoważnym językiem opisu granic.
W każdym wyniku jawnie zaznaczymy konieczne założenia.

\section{Zbieżność ciągów i jej podstawowe własności}
\label{sec:topog-zbieznosc-ciagi}

\begin{definicja}[Granica ciągu]
\label{def:topog-granica-ciagu}
Niech $(X,\tau)$ będzie przestrzenią topologiczną, a
$(x_n)_{n\ge1}$ ciągiem punktów $X$. Mówimy, że
\emph{ciąg $(x_n)$ jest zbieżny do $x\in X$}, i zapisujemy
$x_n\to x$, jeżeli dla każdego otoczenia $U$ punktu $x$
istnieje $N\ge1$ takie, że $x_n\in U$ dla wszystkich $n\ge N$.
Wystarczy sprawdzać otoczenia otwarte, ponieważ każde otoczenie
zawiera otwarte otoczenie punktu $x$.
\end{definicja}

\begin{stwierdzenie}[Zależność od bazy otoczeń]
\label{prop:topog-ciag-baza}
Niech $\mathcal B_x$ będzie bazą otoczeń punktu $x$ w przestrzeni
$(X,\tau)$. Ciąg $(x_n)$ jest zbieżny do $x$ wtedy i tylko wtedy,
gdy dla każdego $B\in\mathcal B_x$ wszystkie wyrazy $(x_n)$,
poza skończoną liczbą, należą do $B$.
\end{stwierdzenie}
\begin{proof}
Konieczność wynika bezpośrednio z definicji granicy. Odwrotnie,
niech $U$ będzie otoczeniem $x$. Z definicji bazy istnieje
$B\in\mathcal B_x$ zawarty w $U$. Od pewnego indeksu
$x_n\in B\subseteq U$, co dowodzi zbieżności.
\end{proof}

\begin{stwierdzenie}[Kryterium metryczne]
\label{prop:topog-ciag-metryka}
Jeżeli topologia przestrzeni $X$ pochodzi od metryki $d$,
to $x_n\to x$ jest równoważne warunkowi:
\[
\forall\varepsilon>0\ \exists N\ge1\ \forall n\ge N:
\quad d(x_n,x)<\varepsilon.
\]
\end{stwierdzenie}
\begin{proof}
Kule $B_d(x,\varepsilon)$, $\varepsilon>0$, tworzą bazę otoczeń
$x$. Zastosowanie stwierdzenia \ref{prop:topog-ciag-baza}
daje dokładnie zapisany warunek.
\end{proof}

\begin{przyklad}[Ciąg liczb odwrotnych]
\label{ex:topog-ciag-odwrotnosci}
W prostej euklidesowej $x_n=1/n$ jest zbieżny do $0$.
Istotnie, dla $\varepsilon>0$ wybieramy liczbę naturalną
$N>1/\varepsilon$. Wówczas dla $n\ge N$ zachodzi
$|1/n-0|=1/n\le1/N<\varepsilon$.
Zauważmy, że żaden wyraz ciągu nie jest równy granicy.
\end{przyklad}

\begin{figure}[htbp]
\centering
\begin{tikzpicture}[x=1.3cm,y=1cm,>=Stealth]
  \fill[blue!9] (-1.20,-.48) rectangle (1.20,.48);
  \draw[blue!70!black,dashed] (-1.20,-.62)--(-1.20,.65);
  \draw[blue!70!black,dashed] (1.20,-.62)--(1.20,.65);
  \draw[->,thick] (-2,0)--(3.4,0);
  \node[below] at (-1.20,-.58) {$-\varepsilon$};
  \node[below] at (1.20,-.58) {$\varepsilon$};
  \fill (0,0) circle (2pt) node[below left] {$0$};
  \fill[orange!85!black] (3,0) circle (2.5pt) node[above] {$x_1$};
  \fill[orange!85!black] (1.5,0) circle (2.5pt) node[above right] {$x_2$};
  \fill[blue!70!black] (1,0) circle (2.5pt) node[above] {$x_3$};
  \fill[blue!70!black] (.60,0) circle (2.5pt) node[above] {$x_5$};
  \fill[blue!70!black] (.30,0) circle (2.5pt);
  \fill[blue!70!black] (.15,0) circle (2.5pt);
  \node[above,blue!70!black] at (0,.62) {$B(0,\varepsilon)$};
  \node[below,font=\scriptsize] at (.25,-1.05) {wyrazy dalsze};
\end{tikzpicture}
\caption{Schemat zbieżności $x_n=1/n$ do $0$: dla ustalonej
kuli wokół granicy wszystkie dostatecznie późne wyrazy
znajdują się wewnątrz. Na rysunku wybrano
$\varepsilon=2/5$ po przeskalowaniu poziomej osi.}
\label{fig:topog-ciag-kula}
\end{figure}

\begin{stwierdzenie}[Podciągi i obrazy ciągłe]
\label{prop:topog-ciag-podciag-ciaglosc}
Niech $(X,\tau_X)$ i $(Y,\tau_Y)$ będą przestrzeniami topologicznymi.
\begin{enumerate}[label=(\roman*)]
\item Jeżeli $x_n\to x$, a $n_1<n_2<\cdots$,
      to $x_{n_k}\to x$.
\item Jeżeli odwzorowanie $f\colon X\to Y$ jest ciągłe w $x$
      oraz $x_n\to x$, to $f(x_n)\to f(x)$.
\end{enumerate}
\end{stwierdzenie}
\begin{proof}
(i) Każde otoczenie $U$ punktu $x$ zawiera wszystkie $x_n$
dla $n\ge N$. Ponieważ $n_k\ge k$, mamy $n_k\ge N$
dla $k\ge N$, a więc $x_{n_k}\in U$.
(ii) Niech $V$ będzie otwartym otoczeniem $f(x)$. Z ciągłości
w $x$ istnieje otwarte $U\ni x$ takie, że $f(U)\subseteq V$.
Od pewnego indeksu $x_n\in U$, zatem $f(x_n)\in V$.
\end{proof}

\begin{twierdzenie}[Jednoznaczność granicy w przestrzeniach Hausdorffa]
\label{thm:topog-ciag-hausdorff}
Jeżeli $X$ jest przestrzenią Hausdorffa, to każdy ciąg w $X$
ma co najwyżej jedną granicę.
\end{twierdzenie}
\begin{proof}
Załóżmy, że $x_n\to x$ oraz $x_n\to y$, gdzie $x\ne y$.
Istnieją rozłączne otwarte otoczenia $U$ punktu $x$ i $V$
punktu $y$. Dla dostatecznie dużego $n$ wyraz $x_n$ należy
jednocześnie do $U$ i $V$, co przeczy $U\cap V=\varnothing$.
\end{proof}

\begin{przyklad}[Brak jednoznaczności]
\label{ex:topog-ciag-wiele-granic}
W dwupunktowej przestrzeni antydyskretnej
$X=\{a,b\}$ jedynym otwartym otoczeniem dowolnego punktu
jest $X$. Dlatego każdy ciąg, w tym ciąg stały $x_n=a$,
zbiega zarówno do $a$, jak i do $b$. Nie można więc
opuszczać aksjomatu Hausdorffa w
 twierdzeniu \ref{thm:topog-ciag-hausdorff}.
\end{przyklad}

\section{Co mogą, a czego nie mogą wykryć ciągi}
\label{sec:topog-ciagi-ograniczenia}

\begin{definicja}[Domknięcie ciągowe]
\label{def:topog-domkniecie-ciagowe}
Dla podzbioru $A\subseteq X$ definiujemy
\emph{zbiór granic ciągów o wyrazach w $A$}:
\[
\operatorname{scl}(A)=
\{x\in X:\ \exists(x_n)_{n\ge1}\in A^{\mathbb N},\ x_n\to x\}.
\]
Zbiór $A$ nazywamy \emph{ciągowo domkniętym}, gdy
$\operatorname{scl}(A)=A$.
\end{definicja}

\begin{stwierdzenie}[Relacja z domknięciem]
\label{prop:topog-ciagi-domkniecie}
Dla każdego $A\subseteq X$ zachodzi
$A\subseteq\operatorname{scl}(A)\subseteq\overline A$.
W szczególności każdy zbiór domknięty jest ciągowo domknięty.
\end{stwierdzenie}
\begin{proof}
Jeżeli $a\in A$, to ciąg stały $x_n=a$ zbiega do $a$,
stąd pierwsze zawieranie. Jeżeli $x_n\in A$ i $x_n\to x$,
to każde otoczenie $U$ punktu $x$ zawiera pewne $x_n$,
a więc $U\cap A\ne\varnothing$. Z kryterium domknięcia
wynika $x\in\overline A$. Jeśli $A$ jest domknięty,
$\overline A=A$, więc oba zawierania stają się równościami.
\end{proof}

\begin{definicja}[Pierwszy aksjomat przeliczalności]
\label{def:topog-pierwsza-przeliczalnosc}
Mówimy, że przestrzeń topologiczna $X$ \emph{spełnia pierwszy
aksjomat przeliczalności}, jeżeli każdy punkt $x\in X$
ma co najwyżej przeliczalną bazę otoczeń.
Warunek jest automatycznie spełniony dla $X=\varnothing$.
\end{definicja}

\begin{przyklad}[Przestrzenie metryczne]
\label{ex:topog-metryczna-pierwsza-przeliczalnosc}
W każdej przestrzeni metrycznej baza
$\{B_d(x,1/n):n\ge1\}$ jest przeliczalna.
Jeżeli $U$ jest otwartym otoczeniem $x$, istnieje $r>0$
z $B_d(x,r)\subseteq U$; wybór $n$ takiego, że $1/n<r$,
daje $B_d(x,1/n)\subseteq U$. Każda przestrzeń metryczna
spełnia więc pierwszy aksjomat przeliczalności.
\end{przyklad}

\begin{twierdzenie}[Ciągi wykrywają domknięcie przy pierwszym aksjomacie przeliczalności]
\label{thm:topog-pierwsza-przeliczalnosc-domkniecie}
Niech $X$ będzie przestrzenią topologiczną spełniającą
pierwszy aksjomat przeliczalności,
a $A\subseteq X$. Wówczas
\[
 x\in\overline A
 \quad\Longleftrightarrow\quad
 \text{istnieje ciąg }(a_n)_{n\ge1}\text{ w }A
 \text{ taki, że }a_n\to x.
\]
W szczególności $A$ jest domknięty wtedy i tylko wtedy,
gdy jest ciągowo domknięty.
\end{twierdzenie}
\begin{proof}
Implikację od prawej do lewej wykazaliśmy w
stwierdzeniu \ref{prop:topog-ciagi-domkniecie}.
Niech teraz $x\in\overline A$ i niech
$\{B_n:n\ge1\}$ będzie bazą otoczeń $x$
(w przypadku bazy skończonej numerujemy jej elementy
z powtórzeniami). Zastępując $B_n$ przez $W_n=B_1\cap\cdots\cap B_n$,
otrzymujemy malejącą bazę otoczeń punktu $x$.
Każdy $W_n$ jest otoczeniem $x$ i przecina $A$,
ponieważ $x\in\overline A$.
Wybieramy $a_n\in A\cap W_n$. Dla dowolnego otoczenia
$U$ punktu $x$ istnieje $k$ z $B_k\subseteq U$.
Gdy $n\ge k$, mamy $a_n\in W_n\subseteq B_k\subseteq U$.
Zatem $a_n\to x$.
Wreszcie, zbiór $A$ jest domknięty dokładnie wtedy,
gdy $\overline A=A$, a na mocy dowiedzionej równoważności
$\overline A=\operatorname{scl}(A)$.
\end{proof}

\begin{twierdzenie}[Ciągowe kryterium ciągłości]
\label{thm:topog-ciagowe-kryterium-ciaglosci}
Niech $X$ będzie przestrzenią spełniającą
pierwszy aksjomat przeliczalności,
$Y$ dowolną przestrzenią topologiczną, $f\colon X\to Y$
i $x\in X$. Odwzorowanie $f$ jest ciągłe w $x$
wtedy i tylko wtedy, gdy dla każdego ciągu $x_n\to x$
zachodzi $f(x_n)\to f(x)$.
\end{twierdzenie}
\begin{proof}
Implikacja wprost wynika ze stwierdzenia
\ref{prop:topog-ciag-podciag-ciaglosc}.
Załóżmy warunek ciągowy i przypuśćmy, że $f$
nie jest ciągłe w $x$. Istnieje otwarte otoczenie
$V$ punktu $f(x)$ takie, że żadne otoczenie $U$
punktu $x$ nie jest zawarte w $f^{-1}(V)$.
Niech $W_n$ będą malejącą bazą otoczeń $x$,
jak w dowodzie twierdzenia
\ref{thm:topog-pierwsza-przeliczalnosc-domkniecie}.
Dla każdego $n$ możemy wybrać
$x_n\in W_n\setminus f^{-1}(V)$.
Wówczas $x_n\to x$, lecz wszystkie $f(x_n)$ leżą poza
$V$, zatem $f(x_n)$ nie może zbiegać do $f(x)$.
Sprzeczność dowodzi ciągłości $f$ w $x$.
\end{proof}

\begin{przyklad}[Topologia koprzeliczalna: ciągi nie wystarczają]
\label{ex:topog-ciagi-koprzeliczalna}
Niech $X=\R$ z topologią koprzeliczalną:
zbiór jest otwarty wtedy i tylko wtedy, gdy jest pusty
lub jego dopełnienie jest co najwyżej przeliczalne.
\begin{enumerate}[label=(\roman*)]
\item Jeżeli $x_n\to x$, to $(x_n)$ jest od pewnego miejsca
stały, równy $x$. Istotnie, zbiór
$C=\{x_n:n\ge1,\ x_n\ne x\}$ jest przeliczalny i nie
zawiera $x$. Otoczenie otwarte $X\setminus C$ punktu $x$
musi zawierać wszystkie wyrazy dostatecznie późne,
a jedyną możliwą ich wartością jest $x$.
Odwrotnie, każdy ciąg ostatecznie stały równy $x$
zbiega do $x$.
\item Każdy podzbiór $A\subseteq X$ jest ciągowo domknięty:
każdy ciąg o wyrazach w $A$ i granicy $x$ jest ostatecznie
stały równy $x\in A$.
Jednak nie każdy podzbiór jest domknięty. Dla przykładu
$A=(0,\infty)$ ma domknięcie $X$, ponieważ przecina
każdy niepusty zbiór otwarty (ten ma przeliczalne dopełnienie,
a $A$ jest nieprzeliczalny).
\item Identyczność $f\colon(X,\tau_{\mathrm{koprz}})
\to(X,\tau_{\mathrm{dysk}})$ zachowuje granice wszystkich
ciągów, bo są one ostatecznie stałe. Nie jest jednak
ciągła: przeciwobraz jednoelementowego zbioru otwartego
w topologii dyskretnej nie jest otwarty w topologii
koprzeliczalnej.
\item Granica dowolnego zbieżnego ciągu jest jednoznaczna,
chociaż $X$ nie jest przestrzenią Hausdorffa:
dwa niepuste zbiory otwarte zawsze się przecinają,
gdyż suma dwóch zbiorów przeliczalnych nie wyczerpuje $X$.
\end{enumerate}
Przestrzeń ta nie spełnia pierwszego aksjomatu przeliczalności:
gdyby była, z twierdzenia
\ref{thm:topog-pierwsza-przeliczalnosc-domkniecie}
wynikałoby, że ciągowo domknięty zbiór $A$ jest domknięty,
wbrew punktowi (ii).
Przykład pokazuje więc, że bez pierwszego aksjomatu przeliczalności
ciągi nie dają ogólnego kryterium domknięcia ani ciągłości.
\end{przyklad}

\section{Sieci i zbieżność Moore'a--Smitha}
\label{sec:topog-sieci}

\begin{definicja}[Zbiór skierowany]
\label{def:topog-zbior-skierowany}
\emph{Zbiorem skierowanym} nazywamy niepusty zbiór $D$
wyposażony w relację $\preceq$ zwrotną, przechodnią
i spełniającą warunek: dla każdych $d,e\in D$ istnieje
$c\in D$ takie, że $d\preceq c$ oraz $e\preceq c$.
Nie wymagamy antysymetrii relacji.
\end{definicja}

\begin{definicja}[Sieć i jej granica]
\label{def:topog-siec-granica}
Niech $(D,\preceq)$ będzie zbiorem skierowanym,
a $X$ przestrzenią topologiczną. \emph{Siecią w $X$}, czyli \emph{ciągiem uogólnionym},
nazywamy rodzinę $(x_d)_{d\in D}$ punktów $X$.
Sieć jest \emph{zbieżna do $x\in X$}, gdy dla każdego
otoczenia $U$ punktu $x$ istnieje $d_0\in D$ takie,
że $x_d\in U$ dla wszystkich $d\succeq d_0$.
Zapisujemy wtedy $x_d\to x$.
\end{definicja}

\begin{przyklad}[Ciąg jako szczególna sieć]
\label{ex:topog-ciag-jako-siec}
Zbiór $\mathbb N=\{1,2,\ldots\}$ z naturalnym porządkiem
jest skierowany. Definicja granicy sieci o indeksach
w $\mathbb N$ jest dokładnie definicją granicy ciągu.
\end{przyklad}

\begin{przyklad}[Indeksowanie skończonymi zbiorami]
\label{ex:topog-siec-skonczone}
Niech $I$ będzie dowolnym niepustym zbiorem.
Rodzina $\mathcal D=\{F\subseteq I:F\text{ jest skończony}\}$
z porządkiem przez zawieranie $\subseteq$ jest skierowana:
dla $F,G$ wspólnym ograniczeniem górnym jest $F\cup G$.
Dzięki temu sieci można indeksować skończonymi
zbiorami nawet wtedy, gdy $I$ jest nieprzeliczalny.
\end{przyklad}

\begin{figure}[htbp]
\centering
\begin{tikzpicture}[>=Stealth, node distance=1.3cm,
  every node/.style={draw,rounded corners=2pt,fill=blue!5,inner sep=5pt}]
  \node (empty) at (0,0) {$\varnothing$};
  \node (a) at (-2,1.25) {$\{a\}$};
  \node (b) at (2,1.25) {$\{b\}$};
  \node (ab) at (0,2.6) {$\{a,b\}$};
  \node (abc) at (0,3.9) {$\{a,b,c\}$};
  \draw[->] (empty)--(a);
  \draw[->] (empty)--(b);
  \draw[->] (a)--(ab);
  \draw[->] (b)--(ab);
  \draw[->] (ab)--(abc);
\end{tikzpicture}
\caption{Fragment zbioru skierowanego skończonych podzbiorów
$I$: dwa indeksy $\{a\}$ i $\{b\}$ mają wspólne ograniczenie
górne $\{a,b\}$. Nie przedstawiono wszystkich indeksów
ani wszystkich strzałek relacji zawierania.}
\label{fig:topog-skierowanie}
\end{figure}

\begin{stwierdzenie}[Obraz sieci zbieżnej]
\label{prop:topog-siec-ciaglosc-jeden-kierunek}
Jeżeli $f\colon X\to Y$ jest ciągłe w $x\in X$
oraz sieć $(x_d)_{d\in D}$ zbiega do $x$, to
$f(x_d)\to f(x)$ w $Y$.
\end{stwierdzenie}
\begin{proof}
Dla dowolnego otoczenia $V$ punktu $f(x)$
z ciągłości $f$ w $x$ wynika, że $f^{-1}(V)$ jest
otoczeniem $x$. Stąd istnieje $d_0$ takie, że
dla $d\succeq d_0$ zachodzi
$x_d\in f^{-1}(V)$, czyli $f(x_d)\in V$.
\end{proof}

\begin{twierdzenie}[Domknięcie jest wykrywane przez sieci]
\label{thm:topog-siec-domkniecie}
Niech $X$ będzie dowolną przestrzenią topologiczną,
$A\subseteq X$, $x\in X$. Wtedy
\[
 x\in\overline A
 \quad\Longleftrightarrow\quad
 \text{istnieje sieć $(a_d)_{d\in D}$ o wyrazach w $A$,
 zbieżna do $x$}.
\]
W szczególności $A$ jest domknięty wtedy i tylko wtedy,
gdy granica każdej zbieżnej sieci o wyrazach w $A$
należy do $A$.
\end{twierdzenie}
\begin{proof}
Jeżeli sieć $(a_d)$ o wyrazach w $A$ zbiega do $x$,
to każde otoczenie $U$ punktu $x$ zawiera jakiś
$a_d\in A$, więc $x\in\overline A$.

Odwrotnie, załóżmy $x\in\overline A$. Utwórzmy
\[
 D=\{(U,a):\ U\text{ jest otwartym otoczeniem }x,
                         \ a\in U\cap A\}.
\]
Zbiór $D$ jest niepusty: $X$ jest otwartym otoczeniem
$x$ i $A\ne\varnothing$ (ponieważ $x\in\overline A$).
Określamy $(U,a)\preceq(V,b)$, gdy $V\subseteq U$.
Relacja jest zwrotna i przechodnia. Jeśli
$(U,a),(V,b)\in D$, to otoczenie $U\cap V$ przecina $A$;
wybieramy $c\in U\cap V\cap A$ i otrzymujemy
$(U\cap V,c)$, które jest ograniczeniem górnym obu indeksów.
Zatem $D$ jest skierowany.
Sieć $a_{(U,a)}=a$ ma wartości w $A$.
Dla dowolnego otwartego otoczenia $W$ punktu $x$
wybieramy $a\in W\cap A$ i indeks $d_0=(W,a)$.
Gdy $(V,b)\succeq d_0$, mamy $V\subseteq W$
i $a_{(V,b)}=b\in V\subseteq W$. Sieć zbiega do $x$.

Ostatnia równoważność wynika z wykazanej charakterystyki:
$A$ jest domknięty dokładnie wtedy, gdy
$\overline A=A$.
\end{proof}

\begin{uwaga}[Unikanie niejawnego wyboru]
\label{rem:topog-siec-bez-wyboru}
Indeksy $(U,a)$ obejmują \emph{wszystkie} dopuszczalne
pary, dlatego do skonstruowania sieci nie musimy
jednocześnie wybrać po jednym punkcie z każdej rodziny
$U\cap A$. Podobne indeksowanie wykorzystamy w dowodzie
jednoznaczności granic.
\end{uwaga}

\begin{twierdzenie}[Ciągłość przez zachowanie granic sieci]
\label{thm:topog-siec-kryterium-ciaglosci}
Niech $X,Y$ będą dowolnymi przestrzeniami topologicznymi,
$f\colon X\to Y$ oraz $x\in X$. Wtedy $f$ jest ciągłe w $x$
dokładnie wtedy, gdy dla każdej sieci $(x_d)$ zbieżnej
do $x$ sieć $(f(x_d))$ zbiega do $f(x)$.
\end{twierdzenie}
\begin{proof}
Konieczność wynika ze stwierdzenia
\ref{prop:topog-siec-ciaglosc-jeden-kierunek}.
Jeśli $f$ nie jest ciągłe w $x$, istnieje otwarty
zbiór $V\ni f(x)$ taki, że $f^{-1}(V)$ nie jest
otoczeniem $x$. Wtedy każde otwarte otoczenie $U$
punktu $x$ przecina zbiór
$A=X\setminus f^{-1}(V)$, czyli $x\in\overline A$.
Z twierdzenia \ref{thm:topog-siec-domkniecie} istnieje
sieć $(a_d)$ o wyrazach w $A$ zbieżna do $x$.
Jednak $f(a_d)\notin V$ dla każdego $d$, więc sieć
$f(a_d)$ nie zbiega do $f(x)$: przeczy temu otoczenie $V$.
Dowodzi to implikacji odwrotnej.
\end{proof}

\begin{twierdzenie}[Jednoznaczność granic sieci i aksjomat Hausdorffa]
\label{thm:topog-sieci-hausdorff}
Dla przestrzeni topologicznej $X$ równoważne są:
\begin{enumerate}[label=(\roman*)]
\item $X$ jest przestrzenią Hausdorffa;
\item każda sieć w $X$ ma co najwyżej jedną granicę.
\end{enumerate}
\end{twierdzenie}
\begin{proof}
Załóżmy (i) i niech $(x_d)$ zbiega do różnych
punktów $x,y\in X$. Wybieramy rozłączne otwarte
otoczenia $U\ni x$ i $V\ni y$. Istnieją indeksy
$d_U,d_V$ takie, że odpowiednie ogony sieci są
zawarte w $U$ i $V$. Z kierunkowości istnieje
$d_0\succeq d_U,d_V$, a wtedy $x_{d_0}\in U\cap V$,
co jest sprzeczne z rozłącznością.

Niech teraz $X$ nie będzie Hausdorffa. Istnieją różne
$x,y\in X$ takie, że każde otwarte otoczenia
$U\ni x$ i $V\ni y$ przecinają się.
Rozważmy wszystkie trójki
\[
 D=\{(U,V,z):\ U\ni x,\ V\ni y\text{ są otwarte},
                                 \ z\in U\cap V\}.
\]
Porządek określamy przez
$(U,V,z)\preceq(U',V',z')$, jeśli
$U'\subseteq U$ i $V'\subseteq V$.
Indeksy mają wspólne ograniczenia górne: dla dwóch
trójek bierzemy otoczenia $U\cap U'$ i $V\cap V'$
oraz dowolny punkt ich niepustego przecięcia.
Sieć $z_{(U,V,z)}=z$ zbiega zarówno do $x$, jak i $y$.
Istotnie, dla otwartego otoczenia $W\ni x$
możemy rozpocząć od trójki $(W,X,z)$,
a wszystkie jej późniejsze indeksy mają pierwszą
współrzędną zawartą w $W$, więc odpowiadające im
punkty należą do $W$. Analogicznie dla $y$,
zaczynając od $(X,W',z)$.
Zatem (ii) nie zachodzi, co kończy dowód.
\end{proof}

\begin{uwaga}
W przeciwieństwie do sieci, jednoznaczność granic
\emph{samych ciągów} nie charakteryzuje aksjomatu
Hausdorffa: kontrprzykład daje topologia koprzeliczalna
z przykładu \ref{ex:topog-ciagi-koprzeliczalna}.
\end{uwaga}

\pagebreak[4]
\begin{przyklad}[Sieć zbieżna, której nie można zastąpić ciągiem]
\label{ex:topog-siec-koprzeliczalna-bez-ciagu}
Rozważmy $X=\R$ z topologią koprzeliczalną,
$A=(0,\infty)$ i punkt $0\notin A$.
Jak pokazaliśmy w przykładzie
\ref{ex:topog-ciagi-koprzeliczalna},
$0\in\overline A$, lecz żaden ciąg o wyrazach w $A$
nie zbiega do $0$: każdy ciąg zbieżny do $0$ musiałby
być od pewnego miejsca stale równy $0$.
Skonstruujemy natomiast konkretną sieć w $A$ zbieżną do $0$.

Niech $D$ będzie zbiorem wszystkich par $(C,a)$, w których
$C\subseteq A$ jest co najwyżej przeliczalny,
a $a\in A\setminus C$. Określamy
$(C,a)\preceq(E,b)$, gdy $C\subseteq E$.
Zbiór $D$ jest niepusty, gdyż $(\varnothing,a)\in D$
dla każdego $a\in A$. Relacja jest zwrotna i przechodnia.
Dla dowolnych $(C,a),(E,b)\in D$ suma $C\cup E$
jest przeliczalna, a $A$ jest nieprzeliczalny,
więc istnieje $c\in A\setminus(C\cup E)$.
Wówczas $(C\cup E,c)$ jest wspólnym ograniczeniem
górnym tych indeksów. Zatem $D$ jest skierowany.
Nie wymagamy tu antysymetrii relacji (zob.
definicję \ref{def:topog-zbior-skierowany}).

Niech $x_{(C,a)}=a$. Zbiór wartości tej sieci jest
dokładnie równy $A$, gdyż $(\varnothing,a)\in D$ dla
każdego $a\in A$. Dla dowolnego otwartego otoczenia
$U$ punktu $0$ zbiór $C=A\setminus U$ jest co najwyżej
przeliczalny. Wybieramy $a\in A\setminus C$
i indeks $d_0=(C,a)$. Gdy $(E,b)\succeq d_0$,
to $C\subseteq E$, zatem
$x_{(E,b)}=b\in A\setminus E\subseteq A\setminus C\subseteq U$.
To dowodzi $x_d\to0$, mimo że żadna z wartości sieci
nie jest równa $0$ i żaden ciąg w $A$ nie zbiega do $0$.
Konstrukcja nie wymaga wyboru jednego punktu dla każdego
zbioru $C$: wszystkie dopuszczalne pary stanowią indeksy.
\end{przyklad}

\section{Podsieci, punkty skupienia sieci i iloczyny}
\label{sec:topog-podsieci}

\begin{definicja}[Podsieć]
\label{def:topog-podsiec}
Niech $(x_d)_{d\in D}$ będzie siecią, a $E$ innym
zbiorem skierowanym. Sieć $(x_{\phi(e)})_{e\in E}$
nazywamy \emph{podsiecią} $(x_d)$, jeśli
$\phi\colon E\to D$ spełnia warunek:
dla każdego $d_0\in D$ istnieje $e_0\in E$ takie,
że $\phi(e)\succeq d_0$ dla wszystkich $e\succeq e_0$.
Nie wymagamy, aby $\phi$ była iniekcją ani
monotoniczna; wymagamy jedynie wskazanej ostatecznej
współkońcowości.
\end{definicja}

\begin{stwierdzenie}[Granice podsieci]
\label{prop:topog-podsiec-granica}
Każda podsieć sieci zbieżnej do $x$ jest zbieżna do $x$.
\end{stwierdzenie}
\begin{proof}
Niech $U$ będzie otoczeniem $x$ i niech
$x_d\in U$ dla wszystkich $d\succeq d_0$.
Z warunku na $\phi$ istnieje $e_0$ takie,
że $\phi(e)\succeq d_0$ dla wszystkich $e\succeq e_0$.
Stąd $x_{\phi(e)}\in U$ dla wszystkich $e\succeq e_0$.
\end{proof}

\begin{definicja}[Punkt skupienia sieci]
\label{def:topog-punkt-skupienia-sieci}
Punkt $x\in X$ nazywamy \emph{punktem skupienia sieci}
$(x_d)_{d\in D}$, jeśli dla każdego otoczenia $U$
punktu $x$ i każdego $d_0\in D$ istnieje
$d\succeq d_0$ takie, że $x_d\in U$.
\end{definicja}

\begin{twierdzenie}[Punkt skupienia a zbieżna podsieć]
\label{thm:topog-punkt-skupienia-podsiec}
Niech $(x_d)_{d\in D}$ będzie siecią w przestrzeni
$X$, a $x\in X$. Wtedy $x$ jest punktem skupienia
sieci dokładnie wtedy, gdy istnieje jej podsieć
zbieżna do $x$.
\end{twierdzenie}
\begin{proof}
Jeżeli podsieć $(x_{\phi(e)})$ zbiega do $x$,
to dla otoczenia $U$ istnieje indeks $e_U$,
od którego wszystkie jej wyrazy leżą w $U$.
Dla zadanego $d_0\in D$ istnieje też $e_d$,
od którego $\phi(e)\succeq d_0$.
Wybierając wspólne ograniczenie górne dla $e_U,e_d$,
otrzymujemy $d=\phi(e)\succeq d_0$ z $x_d\in U$.

Przypuśćmy odwrotnie, że $x$ jest punktem skupienia.
Niech
\[
 E=\{(d,U):d\in D,\ U\text{ jest otwartym otoczeniem }x,
                                        \ x_d\in U\}.
\]
Zbiór $E$ jest niepusty, gdyż dla $U=X$
i dowolnego $d\in D$ mamy $(d,X)\in E$.
Określmy $(d,U)\preceq(e,V)$, gdy $d\preceq e$
oraz $V\subseteq U$. Jest to relacja zwrotna i przechodnia.
Dla $(d,U),(e,V)\in E$ bierzemy $d_0\succeq d,e$.
Z własności punktu skupienia istnieje $c\succeq d_0$
takie, że $x_c\in U\cap V$.
Indeks $(c,U\cap V)$ jest wspólnym ograniczeniem górnym,
a więc $E$ jest skierowany.
Definiujemy $\phi(d,U)=d$. Dla każdego $d_0\in D$
wybieramy $c\succeq d_0$ i zauważamy, że
$(c,X)\in E$. Dla wszystkich późniejszych
$(d,U)\in E$ zachodzi $d\succeq c\succeq d_0$.
Zatem $\phi$ spełnia warunek podsieci.
Wreszcie, dla otwartego otoczenia $W\ni x$
wybieramy $d$ z $x_d\in W$ (istnieje z własności
punktu skupienia) i indeks $(d,W)\in E$.
Dla każdego późniejszego $(e,V)$ mamy $V\subseteq W$,
a ponieważ $x_e\in V$, zachodzi $x_e\in W$.
Podsieć $(x_{\phi(d,U)})$ zbiega do $x$.
\end{proof}

\begin{twierdzenie}[Zbieżność w iloczynie]
\label{thm:topog-siec-iloczyn}
Niech $\{X_i:i\in I\}$ będzie rodziną przestrzeni
topologicznych, a $X=\prod_{i\in I}X_i$ jej iloczynem
z topologią produktową. Niech $(x_d)_{d\in D}$
będzie siecią w $X$, a $x\in X$. Wówczas
\[
 x_d\to x\text{ w }X
 \quad\Longleftrightarrow\quad
 \pi_i(x_d)\to\pi_i(x)\text{ w }X_i
 \text{ dla każdego }i\in I.
\]
\end{twierdzenie}
\begin{proof}
Implikacja wprost wynika z ciągłości rzutowań
$\pi_i\colon X\to X_i$, wykazanej przy konstrukcji
 topologii produktowej.
Odwrotnie, niech wszystkie sieci współrzędnych
zbiegają do właściwych granic. Wystarczy sprawdzić
zbieżność na bazie otoczeń punktu $x$.
Jej elementami są zbiory postaci
$B=\bigcap_{i\in F}\pi_i^{-1}(U_i)$,
gdzie $F\subseteq I$ jest skończony, a $U_i$
jest otwartym otoczeniem $\pi_i(x)$.
Dla każdego $i\in F$ istnieje $d_i\in D$ takie,
że $\pi_i(x_d)\in U_i$ dla wszystkich $d\succeq d_i$.
Z kierunkowości $D$ i skończoności $F$ istnieje
wspólne ograniczenie górne $d_0$ tych indeksów
(dla $F=\varnothing$ wystarczy dowolny $d_0\in D$).
Wtedy dla $d\succeq d_0$ zachodzi
$x_d\in B$. To dowodzi zbieżności w iloczynie.
\end{proof}

\section{Filtry i ich zbieżność}
\label{sec:topog-filtry}

\begin{definicja}[Filtr właściwy]
\label{def:topog-filtr}
Niech $X$ będzie zbiorem. \emph{Filtrem (właściwym)}
na $X$ nazywamy rodzinę $\mathcal F\subseteq\mathcal P(X)$
spełniającą następujące warunki:
\begin{enumerate}[label=(\roman*)]
\item $X\in\mathcal F$ i $\varnothing\notin\mathcal F$;
\item jeżeli $F,G\in\mathcal F$, to $F\cap G\in\mathcal F$;
\item jeśli $F\in\mathcal F$ i $F\subseteq B\subseteq X$,
      to $B\in\mathcal F$.
\end{enumerate}
\emph{Bazą filtru} nazywamy niepustą rodzinę
$\mathcal B\subseteq\mathcal P(X)\setminus\{\varnothing\}$
taką, że dla $B,C\in\mathcal B$ istnieje
$D\in\mathcal B$ z $D\subseteq B\cap C$.
Generuje ona filtr $\mathcal F_{\mathcal B}$ wszystkich
zbiorów zawierających co najmniej jeden element
$\mathcal B$.
\end{definicja}

\begin{proof}[Uzasadnienie konstrukcji filtru z bazy]
Rodzina $\mathcal F_{\mathcal B}$ zawiera $X$,
a nie zawiera $\varnothing$, ponieważ elementy bazy
są niepuste. Jeśli $F$ zawiera $B\in\mathcal B$, a $G$
zawiera $C\in\mathcal B$, to $F\cap G$ zawiera
odpowiedni $D\subseteq B\cap C$, więc należy
do $\mathcal F_{\mathcal B}$. Własność rozszerzania
zbiorów jest wprost wpisana w definicję.
\end{proof}

\begin{definicja}[Granica filtru]
\label{def:topog-filtr-granica}
Niech $X$ będzie przestrzenią topologiczną,
$\mathcal F$ filtrem na $X$ i $x\in X$.
Mówimy, że $\mathcal F$ \emph{zbiega do $x$}, gdy
każde otoczenie punktu $x$ należy do $\mathcal F$.
Zapisujemy $\mathcal F\to x$.
\end{definicja}

\begin{przyklad}[Filtr otoczeń]
\label{ex:topog-filtr-otoczen}
Dla każdego $x\in X$ rodzina wszystkich otoczeń
punktu $x$ jest filtrem na $X$: jest niepusta,
nie zawiera $\varnothing$, jest zamknięta na skończone
przecięcia i na rozszerzanie zbiorów.
Filtr ten zbiega do $x$ z samej definicji.
\end{przyklad}

\begin{twierdzenie}[Domknięcie przez filtry]
\label{thm:topog-filtr-domkniecie}
Dla dowolnej przestrzeni topologicznej $X$,
zbioru $A\subseteq X$ i punktu $x\in X$ warunek
$x\in\overline A$ jest równoważny istnieniu
filtru $\mathcal F$ na $X$ takiego, że
$A\in\mathcal F$ i $\mathcal F\to x$.
\end{twierdzenie}
\begin{proof}
Załóżmy $x\in\overline A$. Niech $\mathcal B$
będzie rodziną zbiorów $A\cap U$, gdzie $U$
przebiega wszystkie otwarte otoczenia $x$.
Każdy element $\mathcal B$ jest niepusty z
charakterystyki domknięcia. Dla $A\cap U$ i $A\cap V$
element $A\cap(U\cap V)$ należy do $\mathcal B$
i jest ich przecięciem. Zatem $\mathcal B$ jest bazą
filtru $\mathcal F$. Ponieważ $A=A\cap X\in\mathcal B$,
mamy $A\in\mathcal F$. Każde otwarte otoczenie
$U$ zawiera $A\cap U\in\mathcal B$, więc należy
do $\mathcal F$; dowolne otoczenie $x$ zawiera
otwarte otoczenie, zatem również należy do filtru.

Odwrotnie, jeżeli $A\in\mathcal F$ i
$\mathcal F\to x$, to dla każdego otwartego
otoczenia $U$ punktu $x$ mamy $U\in\mathcal F$.
Stąd $A\cap U\in\mathcal F$, a ponieważ filtr
jest właściwy, $A\cap U\ne\varnothing$.
Oznacza to $x\in\overline A$.
\end{proof}

\begin{twierdzenie}[Sieci i filtry opisują tę samą zbieżność]
\label{thm:topog-sieci-filtry}
Niech $X$ będzie przestrzenią topologiczną.
\begin{enumerate}[label=(\roman*)]
\item Dla każdej sieci $(x_d)_{d\in D}$ w $X$
rodzina wszystkich zbiorów zawierających
pewien ogon $\{x_d:d\succeq d_0\}$
jest filtrem $\mathcal F_{(x_d)}$.
Sieć zbiega do $x$ wtedy i tylko wtedy,
gdy $\mathcal F_{(x_d)}\to x$.
\item Dla każdego filtru $\mathcal F$ na $X$
istnieje sieć $(x_d)_{d\in D}$ taka, że dla każdego
$x\in X$ zachodzi
$\mathcal F\to x\iff x_d\to x$.
\end{enumerate}
\end{twierdzenie}
\begin{proof}
(i) Ogony są niepuste, bo dla każdego $d_0$
zawierają $x_{d_0}$. Przecięcie dwóch ogonów
zawiera ogon o indeksie będącym wspólnym
ograniczeniem górnym ich indeksów początkowych.
Stanowią zatem bazę filtru. Otoczenie $U$ należy
do generowanego filtru dokładnie wtedy, gdy zawiera
pewien ogon sieci. Jest to definicja zbieżności sieci.

(ii) Niech
\[
 D=\{(F,a):F\in\mathcal F,\ a\in F\}.
\]
Rodzina $D$ jest niepusta, ponieważ $X\in\mathcal F$
i $X\ne\varnothing$ (właściwy filtr na zbiorze pustym
nie istnieje). Porządkujemy ją przez
$(F,a)\preceq(G,b)\iff G\subseteq F$.
Dla dwóch indeksów wybieramy dowolny punkt
$c\in F\cap G$, który istnieje, gdyż
$F\cap G\in\mathcal F$. Wtedy $(F\cap G,c)$
jest wspólnym ograniczeniem górnym.
Definiujemy sieć $x_{(F,a)}=a$.

Jeśli $\mathcal F\to x$, to dla każdego otwartego
otoczenia $U$ punktu $x$ mamy $U\in\mathcal F$.
Wybieramy $a\in U$ i indeks $(U,a)$. Każdy późniejszy
$(G,b)$ spełnia $G\subseteq U$, więc $x_{(G,b)}=b\in U$.
Zatem sieć zbiega do $x$.

Na odwrót, niech $(x_{(F,a)})$ zbiega do $x$.
Dla każdego otwartego otoczenia $U$ punktu $x$
istnieje $(F,a)$ takie, że
$x_{(G,b)}\in U$ dla wszystkich $(G,b)\succeq(F,a)$.
Dla dowolnego $b\in F$ indeks $(F,b)$ należy do $D$
i spełnia $(F,b)\succeq(F,a)$, zatem $b\in U$.
Wynika stąd $F\subseteq U$, a więc $U\in\mathcal F$
przez własność rozszerzania zbiorów. Każde inne
otoczenie punktu $x$ zawiera otwarte, więc również
należy do filtru. Otrzymujemy $\mathcal F\to x$.
\end{proof}

\begin{wniosek}[Ciągłość i jednoznaczność granic filtrów]
\label{cor:topog-filtr-ciaglosc-hausdorff}
Niech $X,Y$ będą przestrzeniami topologicznymi.
\begin{enumerate}[label=(\roman*)]
\item Odwzorowanie $f\colon X\to Y$ jest ciągłe
w $x\in X$ dokładnie wtedy, gdy dla każdego filtru
$\mathcal F\to x$ filtr generowany przez zbiory
$f(F)$, $F\in\mathcal F$, zbiega do $f(x)$.
\item Przestrzeń $X$ jest Hausdorffa dokładnie wtedy,
gdy każdy filtr na $X$ ma co najwyżej jedną granicę.
\end{enumerate}
\end{wniosek}
\begin{proof}
Dla (i) najpierw sprawdźmy, że rodzina
$\{f(F):F\in\mathcal F\}$ jest bazą filtru na $Y$.
Każdy jej element jest niepusty, ponieważ
$\varnothing\notin\mathcal F$. Jeśli $F,G\in\mathcal F$,
to $F\cap G\in\mathcal F$ oraz
$f(F\cap G)\subseteq f(F)\cap f(G)$.
Spełnione są więc warunki bazy filtru z
 definicji \ref{def:topog-filtr}.
Ponadto $V\subseteq Y$ należy do filtru obrazowego
dokładnie wtedy, gdy
$f^{-1}(V)\in\mathcal F$: jeśli $f(F)\subseteq V$
dla pewnego $F\in\mathcal F$, to
$F\subseteq f^{-1}(V)$ i filtr zawiera przeciwobraz;
odwrotnie z $f^{-1}(V)\in\mathcal F$ wynika
$f(f^{-1}(V))\subseteq V$.
Jeżeli $f$ jest ciągłe w $x$ i $\mathcal F\to x$,
to dla każdego otwartego otoczenia $V$ punktu $f(x)$
zbiór $f^{-1}(V)$ jest otoczeniem $x$, zatem należy
do $\mathcal F$. Z wykazanej równoważności
$V$ należy do filtru obrazowego, co dowodzi jego
zbieżności do $f(x)$.

Odwrotnie, przypuśćmy, że $f$ zachowuje granice
wszystkich filtrów, ale nie jest ciągłe w $x$.
Z twierdzenia \ref{thm:topog-siec-kryterium-ciaglosci}
istnieje sieć $(x_d)$ zbieżna do $x$, której obraz
$(f(x_d))$ nie zbiega do $f(x)$. Filtr ogonów
$\mathcal F_{(x_d)}$ zbiega do $x$ na mocy
twierdzenia \ref{thm:topog-sieci-filtry}(i).
Ponieważ ogony sieci obrazowej są dokładnie
obrazami ogonów sieci wyjściowej, ich filtr
jest filtrem obrazowym $f_*(\mathcal F_{(x_d)})$.
Nie zbiega on do $f(x)$, co przeczy założeniu.

Dla (ii) jeśli $X$ jest Hausdorffa, dwa różne punkty
$x,y$ mają rozłączne otwarte otoczenia $U,V$.
Filtr zbieżny do obu musiałby zawierać
$U\cap V=\varnothing$, co jest niemożliwe.
Jeśli $X$ nie jest Hausdorffa, to z dowodu
twierdzenia \ref{thm:topog-sieci-hausdorff}
istnieje sieć zbieżna do dwóch różnych punktów.
Jej filtr ogonów z punktu (i) twierdzenia
\ref{thm:topog-sieci-filtry} ma te same granice.
\end{proof}

\section{Podsumowanie: trzy języki zbieżności}
\label{sec:topog-zbieznosc-zestawienie}

\begin{center}
\begin{tabular}{@{}p{.22\textwidth}p{.31\textwidth}p{.34\textwidth}@{}}
\hline
\textbf{Narzędzie} & \textbf{Co wykrywa zawsze} &
\textbf{Kiedy wystarcza}\ \\ \hline
Ciągi & Granice i konieczny warunek ciągłości &
Dla domknięcia i ciągłości, jeśli przestrzeń
spełnia pierwszy aksjomat przeliczalności.\\
Sieci & Domknięcie, ciągłość, aksjomat Hausdorffa &
W każdej przestrzeni topologicznej.\\
Filtry & Domknięcie, ciągłość, aksjomat Hausdorffa &
W każdej przestrzeni topologicznej.\\
\hline
\end{tabular}
\end{center}

\begin{uwaga}[Powiązania z dalszymi rozdziałami]
\label{rem:topog-zbieznosc-dalsze}
Sieci i filtry pozwolą później sformułować
charakterystyki zwartości oraz zbadać dokładniej
przestrzenie bez przeliczalnej bazy otoczeń.
Warunki przeliczalności i aksjomaty metryzowalności
będą analizowane odrębnie; w szczególności
spełnienia pierwszego aksjomatu przeliczalności nie
utożsamiamy tutaj z metryzowalnością.
\end{uwaga}
